
A partially ordered set (poset) is a set together with a partial order. 
For instance, the set of subsets of a fixed set form an order structure 
with respect to set-inclusion.
The Hasse diagram is a diagram whose nodes represent the element. 
Nodes are arranged from top to bottom, 
and relations are indicated by lines. 
Only the maximal inclusions are drawn. 
The other relations can be inferred from transitivity.

For example, consider the set of subsets of 
a 4-element set, say $\{1,2,3,4\}$:

\orbiteroutput{GRAPHICS/WRAPPERS/subset_lattice_4}

Here, $\{1\}$ is a subset of $\{1,2\}$, which is a subset of $\{1,2,3\}$.
There is no line from $\{1\}$ to $\{1,2,3\}$ as the relation is not maximal. 


\bigskip


Partial orders can be useful in the problem of classifying combinatorial 
objects by means of substructures.
A poset action is group action on a poset. 
Let $G$ be a group and let $\cP$ be a poset. 
For a poset action, we require that 
$$
x \le y \Rightarrow xg \le yg \quad \mbox{for all} \; x,y \in \cP, \; g \in G.
$$
For background on poset actions, see Plesken~\cite{Plesken82}.
The orbits of $G$ on $\cP$ form another poset, the poset of orbits. 
A ranked poset can be decomposed into layers according to the rank function. 
It is then possible to build the poset by considering adjacent levels inductively.
In each step, the next layer is constructed from the previous layer and all the layers below.
This allows one to reduce a classification problem to a sequence of smaller, 
and in many cases easier steps.  


\bigskip




Orbiter uses the algorithm of Schmalz~\cite{Schmalz90} to perform poset classification.
Two versions are available: one for subset-type posets and one for subspace-type posets.
The subspace lattice of $V(3,2)=\bbF_2^3$ is shown below:

\orbiteroutput{GRAPHICS/LATEX/V_3_2}

The basis elements are listed, using the enumerator for 
elements of the projective geometry $\PG(2,2)$ 
explained in Section~\ref{sec:FPG}.




\bigskip


The commands listed in Table~\ref{tab:pc:options}
can be used to control the poset classification algorithm.

\orbitertableofcommands{poset_classification_control}

By default, Orbiter will choose the lexiographically least orbit representatives. 
It is possible to direct Orbiter to choose different orbit representatives. 
To this end, the nodes in the orbit tree are labeled. 
The node number is the zero-based number of a given node in the tree, 
using the breadth first ordering. 

\bigskip

For posets whose orbits have been computed, 
there are activities as shown in Tables~\ref{tab:pc:activity:1}-\ref{tab:pc:activity:2}. 
\orbitertableofcommands{poset_classification_activity_1}
\orbitertableofcommands{poset_classification_activity_2}

\bigskip

For reporting purposes, the options in Table~\ref{tab:pc:report:options} can be used. 
\orbitertableofcommands{poset_classification_report_options}

\bigskip

Suppose that Orbiter chooses element $a$ at node $n.$ 
If we prefer $b$ over $a$, we can issue the command 
$$
\verb'-preferred_choice' \; n \; a \; b
$$
It forces Orbiter to choose $b$ instead of $a$ at node $n$.

